\scanpage{47}
\mspage{4}
\figure{\begin{tikzpicture}[scale=1.4]
\draw (-2.1,0)--(2.2,-.03);\draw (-1,0) arc[start angle=180,end angle=0,radius=1];\draw (0,-1)--(0,3.35);\draw (.5,0)--(.5,3.35);
\draw[thick] (0,1)--(0,3.35) (.5,.866)--(.5,3.35);
\draw (0,1.15)--(.5,1.55) (0,1.7)--(.5,2.1) (0,2.25)--(.5,2.65) (0,2.8)--(.5,3.2) (0,3.15)--(.22,3.35);
\fill (0,1) circle(.025) (.5,.866) circle(.025);\node[left] at (0,1.12){$\uline{i}$};\node[right] at (.5,.866){$\jmath$};
\end{tikzpicture}}
\[
\rho=\deleted{\unsure{\sigma}}\begin{pmatrix}0&1\\-1&1\end{pmatrix}\qquad\rho^3=\omega=\begin{pmatrix}-1&0\\0&-1\end{pmatrix}
\]
\[
\sigma=\begin{pmatrix}0&+1\\-1&0\end{pmatrix}\qquad\sigma^2=\omega\qquad\sigma^{-1}=\omega\sigma
\]
On a choisi les signes de telle façon pour

a) $\rho^3=\deleted{\unsure{1}}-1$ i.e. $\rho$ d'ordre $6$ i.e. il engendre le groupe d'isotropie \inserted{de $\jmath$} (mais cela laisse le choix entre $\rho$ et $\rho^{-1}$)

b) $\sigma$ soit la rotation de l'angle $-\dfrac\pi2$ dans $\mathbb R^2$ (pour la structure euclidienne habituelle) et que $\rho$ soit une rotation d'angle $-2\pi/6$ (pour une métrique invariante convenable).
La chose importante dans cette convention, \unsure{c'est} que pour $\rho$ et $\sigma$, le sens des rotations (i.e. négatif) soit le $\unsure{\overline{\mathrm m}}$.
Quand on interprète \deleted{que} les matrices $\begin{pmatrix}a&b\\c&d\end{pmatrix}\in SL(2,\mathbb Z)$ comme opérant sur des \deleted{\unsure{systèmes}}\inserted{couples} $(z_1,z_2)$ de vecteurs de $\mathbb C$, engendrant un sous-groupe discret de $\mathbb C$ (avec $z_1/z_2\in\mathcal D$), alors l'action de $\rho$ et de $\sigma$ sur $z_1,z_2$ est ``positive'', en ce sens que $(z_1,u(z_1))$ et $(z_2,u(z_2))$ forment des $\mathbb R$-bases \uline{directes} de $\mathbb C\simeq\mathbb R^2$.

Voyons le groupe des relèvements de $SL(2,\mathbb Z)$ opérant sur $P\mathcal D$ sur $\widetilde{P\mathcal D}$, qui est donc une extension centrale de $\mathcal T_{11}=\unsure{GL|SL}(2,\mathbb Z)$ par $\mathbb Z$, \deleted{\unsure{soit}}\inserted{qui peut s'écrire $S\mathcal T_{11}$}\deleted{$\widetilde{SL}(2,\mathbb Z)$}.
Je désigne par \deleted{$\unsure{\omega}$} $\overline\omega$ l'élément de \deleted{$SL(2,\mathbb Z)$} $S\mathcal T_{11}$ défini par le générateur de $\mathbb Z$, opérant sur $\widetilde{P\mathcal D}$ par $(\tau,z)\mathrel{\text{---}}(\tau,z+1)$.
\mspage{46}
